\documentclass[12pt]{article}
\usepackage{amsmath, amssymb}
\usepackage[margin=1in]{geometry}
\setlength{\parskip}{6pt}
\title{QUBO Formulation: Densest K-Subgraph Problem}
\date{}
\begin{document}
\maketitle

\subsection*{Densest K-Subgraph Problem}

\paragraph{Problem Statement.}
Given an undirected weighted graph $G = (V, E)$ with $N = |V|$ vertices and edge weights $w_{ij} \geq 0$ for each $\{i, j\} \in E$, the Densest K-Subgraph Problem asks for a subset of vertices $S \subseteq V$ of cardinality exactly $K$ such that the total weight of the edges induced by $S$ is maximized. Formally, the problem is:
\begin{align}
    \max_{S \subseteq V} \quad & \sum_{\{i, j\} \subseteq S} w_{ij} \\
    \text{s.t.} \quad & |S| = K.
\end{align}

\paragraph{Decision Variables.}
We introduce a binary variable $x_i \in \{0, 1\}$ for each vertex $i \in V$. The variable $x_i$ indicates the selection of vertex $i$:
\begin{align}
    x_i = 
    \begin{cases}
        1 & \text{if } i \in S, \\
        0 & \text{otherwise}.
    \end{cases}
\end{align}
The constraint $|S| = K$ is enforced by $\sum_{i \in V} x_i = K$. The induced edge weight is represented by $\sum_{i < j} w_{ij} x_i x_j$, where the sum ranges over all pairs of distinct vertices.

\paragraph{QUBO Derivation.}
We convert the maximization problem to a minimization problem by negating the objective. The resulting Quadratic Unconstrained Binary Optimization (QUBO) Hamiltonian $H(\mathbf{x})$ is constructed by adding a quadratic penalty term to enforce the cardinality constraint.

\textbf{Step 1: Objective Function.}
The original objective is to maximize $\sum_{i < j} w_{ij} x_i x_j$. For minimization, we define:
\begin{align}
    H_{\text{obj}}(\mathbf{x}) = -\sum_{i < j} w_{ij} x_i x_j.
\end{align}

\textbf{Step 2: Constraint.}
The cardinality constraint is:
\begin{align}
    \sum_{i \in V} x_i = K.
\end{align}

\textbf{Step 3: Penalty Term.}
We introduce a penalty term with weight $\lambda > 0$:
\begin{align}
    H_{\text{pen}}(\mathbf{x}) = \lambda \left( \sum_{i \in V} x_i - K \right)^2.
\end{align}
Expanding the square and using the idempotent property $x_i^2 = x_i$ for binary variables:
\begin{align}
    \left( \sum_{i \in V} x_i - K \right)^2 
    &= \sum_{i \in V} x_i^2 + \sum_{i \neq j} x_i x_j - 2K \sum_{i \in V} x_i + K^2 \\
    &= \sum_{i \in V} x_i + 2 \sum_{i < j} x_i x_j - 2K \sum_{i \in V} x_i + K^2 \\
    &= 2 \sum_{i < j} x_i x_j + (1 - 2K) \sum_{i \in V} x_i + K^2.
\end{align}
Multiplying by $\lambda$, the penalty contribution is:
\begin{align}
    H_{\text{pen}}(\mathbf{x}) = 2\lambda \sum_{i < j} x_i x_j + \lambda(1 - 2K) \sum_{i \in V} x_i + \lambda K^2.
\end{align}

\textbf{Step 4: Combined Hamiltonian.}
Combining the objective and penalty terms, the total QUBO Hamiltonian is:
\begin{align}
    H(\mathbf{x}) = \sum_{i < j} (2\lambda - w_{ij}) x_i x_j + \sum_{i \in V} \lambda(1 - 2K) x_i + \lambda K^2.
\end{align}

\textbf{Step 5: Q Matrix Entries.}
The Hamiltonian is expressed in the standard QUBO form $H(\mathbf{x}) = \mathbf{x}^T Q \mathbf{x} + c$, where $\mathbf{x}^T Q \mathbf{x} = \sum_{i} Q_{ii} x_i + 2 \sum_{i < j} Q_{ij} x_i x_j$. Matching coefficients with $H(\mathbf{x})$ yields:
\begin{align}
    Q_{ij} &= \lambda - \frac{w_{ij}}{2} \quad \text{for } i < j, \\
    Q_{ii} &= \lambda(1 - 2K) \quad \text{for } i \in V, \\
    c &= \lambda K^2.
\end{align}

\paragraph{Penalty Weight Justification.}
The penalty weight $\lambda$ must be chosen sufficiently large to ensure that any violation of the cardinality constraint incurs an energy cost exceeding the maximum possible gain from the objective function. The minimum energy cost for violating the constraint by one unit is $\lambda$. The maximum possible increase in the objective value by altering the selection is bounded by the sum of all edge weights incident to any single vertex, which is at most $N \cdot \max_{i,j} w_{ij}$. Therefore, choosing
\begin{align}
    \lambda > N \cdot \max_{i, j \in V} w_{ij}
\end{align}
guarantees that the penalty cost for any infeasible solution strictly dominates any potential objective improvement, ensuring the global minimum corresponds to a feasible solution.

\paragraph{Correctness Argument.}
Minimizing $H(\mathbf{x})$ over $\mathbf{x} \in \{0, 1\}^N$ recovers the optimal solution to the Densest K-Subgraph Problem. For any feasible solution ($\sum_i x_i = K$), the penalty term vanishes, and $H(\mathbf{x})$ equals the negated objective value. Thus, minimizing $H$ over feasible solutions is equivalent to maximizing the induced edge weight. For any infeasible solution, the penalty term is at least $\lambda$. By the choice of $\lambda$, the energy of any infeasible solution exceeds the energy of the best feasible solution. Consequently, the global minimum of $H(\mathbf{x})$ must be attained at a feasible configuration that maximizes the objective.

\paragraph{Illustrative Example.}
Consider a graph with $N=3$ vertices $\{0, 1, 2\}$ and edge weights $w_{01} = 1$, $w_{02} = 2$, $w_{12} = 3$. Let $K=2$. The optimal solution selects vertices $\{1, 2\}$ with induced weight $3$.

Using the general formulas with $\lambda = 10$ (satisfying $\lambda > 3 \cdot 3 = 9$):
\begin{align}
    Q_{01} &= 10 - \frac{1}{2} = 9.5, \quad Q_{02} = 10 - \frac{2}{2} = 9.0, \quad Q_{12} = 10 - \frac{3}{2} = 8.5, \\
    Q_{00} = Q_{11} = Q_{22} &= 10(1 - 4) = -30, \\
    c &= 10 \cdot 2^2 = 40.
\end{align}
The Hamiltonian is:
\begin{align}
    H(\mathbf{x}) = 9.5 x_0 x_1 + 9.0 x_0 x_2 + 8.5 x_1 x_2 - 30(x_0 + x_1 + x_2) + 40.
\end{align}
Evaluating at the optimal configuration $\mathbf{x}^* = (0, 1, 1)$:
\begin{align}
    H(\mathbf{x}^*) = 8.5(1) - 30(2) + 40 = 8.5 - 60 + 40 = -11.5.
\end{align}
This corresponds to the negated objective value $-3$, confirming correctness. Any infeasible configuration, such as $\mathbf{x} = (1, 1, 0)$, yields $H = 9.5 - 60 + 40 = -10.5$, which is strictly higher than $-11.5$.

\end{document}